\documentclass[tikz,border=0pt]{standalone}
\usepackage[fontset=none]{ctex}
\usepackage{amsmath,amssymb,bm}
\usetikzlibrary{arrows.meta,calc,positioning,angles,quotes}

\definecolor{ink}{HTML}{263238}
\definecolor{muted}{HTML}{6B7780}
\definecolor{blue}{HTML}{245BFF}
\definecolor{green}{HTML}{087F5B}
\definecolor{orange}{HTML}{E67700}
\definecolor{red}{HTML}{E03131}
\definecolor{purple}{HTML}{7048E8}
\definecolor{paleBlue}{HTML}{EDF3FF}
\definecolor{paleGreen}{HTML}{EAF8F3}
\definecolor{paleOrange}{HTML}{FFF4E6}
\definecolor{palePurple}{HTML}{F3F0FF}

\tikzset{
  title/.style={font=\fontsize{18}{22}\selectfont\bfseries,text=ink,align=center},
  subtitle/.style={font=\fontsize{12.5}{16}\selectfont\bfseries,text=ink,align=center},
  body/.style={font=\fontsize{10.5}{14}\selectfont,text=ink,align=center},
  small/.style={font=\fontsize{8.8}{11.5}\selectfont,text=ink,align=center},
  box/.style={draw=muted,line width=.8pt,rounded corners=2pt,inner sep=6pt,body},
  vec/.style={-{Latex[length=3.2mm,width=2.3mm]},line width=1.35pt},
  patharrow/.style={-{Latex[length=2.8mm,width=2mm]},line width=1.25pt},
  nodept/.style={circle,fill=ink,minimum size=5pt,inner sep=0pt}
}

\usepackage{fontspec,xeCJK}
\IfFontExistsTF{Heiti SC}{\setCJKmainfont[AutoFakeBold=1.5]{Heiti SC}}{\setCJKmainfont[AutoFakeBold=1.5]{FandolHei-Regular}}
\IfFontExistsTF{Heiti SC}{\setCJKsansfont{Heiti SC}}{\setCJKsansfont{FandolHei-Regular}}
\usetikzlibrary{bending}
\begin{document}
\begin{tikzpicture}[x=1cm,y=1cm]
  \path[use as bounding box] (0,0) rectangle (16,8); \fill[white] (0,0) rectangle (16,8);
  \node[title] at (8,7.45) {亥姆霍兹线圈：线圈间距等于半径时，中心场最平坦};
  \draw[muted,line width=.8pt] (7.1,1.0)--(7.1,6.75);
  \node[subtitle,blue] at (3.55,6.35) {两只同轴圆线圈，电流方向相同};
  \draw[ink,vec] (.7,3.7)--(6.55,3.7); \node[small,above] at (6.3,3.75) {轴线 $z$};
  \draw[blue,line width=2pt] (2.35,1.7)--(2.35,5.7); \draw[blue,line width=2pt] (4.95,1.7)--(4.95,5.7);
  \draw[red,line width=1.1pt] (2.35,5.25) circle (.14); \fill[red] (2.35,5.25) circle (.04);
  \draw[red,line width=1.1pt] (4.95,5.25) circle (.14); \fill[red] (4.95,5.25) circle (.04);
  \draw[red] (2.35,2.15) circle (.14); \draw[red] (4.95,2.15) circle (.14);
  \draw[red,line width=1pt] (2.25,2.05)--(2.45,2.25); \draw[red,line width=1pt] (2.25,2.25)--(2.45,2.05);
  \draw[red,line width=1pt] (4.85,2.05)--(5.05,2.25); \draw[red,line width=1pt] (4.85,2.25)--(5.05,2.05);
  \draw[green,<->,line width=1.15pt] (2.35,1.3)--(4.95,1.3); \node[small,green,below] at (3.65,1.25) {间距 $d$};
  \draw[orange,<->,line width=1.15pt] (1.95,3.7)--(1.95,5.7); \node[small,orange,left] at (1.85,4.7) {半径 $R$};
  \draw[purple,vec,line width=1.7pt] (3.65,3.7)--(5.8,3.7); \node[small,purple,above] at (4.2,3.75) {$\bm B$};
  \node[box,fill=paleBlue,text width=5.7cm] at (3.55,.95)
    {对称性自动给出 $B'(0)=0$；选择 $d=R$ 还能使 $B''(0)=0$};

  \node[subtitle,green] at (11.55,6.35) {轴线场在中心附近的形状};
  \draw[ink,->] (8.0,1.55)--(15.2,1.55); \draw[ink,->] (11.6,1.2)--(11.6,5.75);
  \node[small,right] at (15.1,1.45) {$z/R$}; \node[small,above] at (11.6,5.65) {$B(z)$};
  \draw[blue,line width=1.7pt,domain=-3.1:3.1,samples=180,variable=\x]
    plot ({11.6+1.0*\x},{1.55+1.25*((1+(\x-.5)^2)^(-1.5)+(1+(\x+.5)^2)^(-1.5))});
  \draw[green,dashed,line width=1.1pt] (10.0,3.34)--(13.2,3.34);
  \node[small,green,above] at (13.25,3.37) {中心附近近似匀强};
  \node[box,fill=paleGreen,text width=6.1cm] at (11.55,.9)
    {$\displaystyle B(0)=\left(\frac45\right)^{3/2}\frac{\mu_0NI}{R}$};
\end{tikzpicture}
\end{document}
